\documentclass[11pt,a4paper]{article}
\input{../../sablony/pripravy/priprava.sty}
\usepackage{tikz,pgfplots}
\pgfplotsset{compat=1.18}
\nastavPripravu{1. ročník}{Kvadratické rovnice a nerovnice}{123}{Kvadratické nerovnice II: metoda intervalů}
\begin{document}
\blok{Kvadratické nerovnice II: metoda intervalů}
Rozklad na součin, kritické body a tabulka znamének. Grafické a algebraické řešení musí souhlasit.
\blok{Metoda a návaznost na graf}
\begin{itemize}
\item Po kořenovém rozkladu určuj znaménko součinu na intervalech; ve výsledku piš sjednocení intervalů.
\item Nezaměňuj podmínku „větší než nula“ s rovnicí rovnou nule.
\end{itemize}
\blok{Společně řešené příklady}
\begin{prikladReseniBox}[title={Příklad 1}]
$x^{2} - 5x + 6\ge0$
\tcblower
\[
\begin{aligned}
f(x)&=x^{2} - 5x + 6\\[2pt]
f(x)=0&\ \Rightarrow\ K_0=\{2;3\}\\[2pt]
K&=(-\infty;2\rangle\cup\langle3;+\infty)
\end{aligned}
\]
\end{prikladReseniBox}
\begin{prikladReseniBox}[title={Příklad 2}]
$- 2x^{2} + 2x + 12>0$
\tcblower
\[
\begin{aligned}
f(x)&=- 2x^{2} + 2x + 12\\[2pt]
f(x)=0&\ \Rightarrow\ K_0=\{-2;3\}\\[2pt]
K&=(-2;3)
\end{aligned}
\]
\end{prikladReseniBox}
\begin{prikladReseniBox}[title={Příklad 3}]
$2x^{2} + 4x - 6<0$
\tcblower
\[
\begin{aligned}
f(x)&=2x^{2} + 4x - 6\\[2pt]
f(x)=0&\ \Rightarrow\ K_0=\{-3;1\}\\[2pt]
K&=(-3;1)
\end{aligned}
\]
\end{prikladReseniBox}
\blok{Samostatné procvičování}
\begin{zadaniReseniBox}[title={Samostatná práce 1}]
$x^{2} + x - 6<0$
\tcblower
\[
\begin{aligned}
f(x)&=x^{2} + x - 6\\[2pt]
f(x)=0&\ \Rightarrow\ K_0=\{-3;2\}\\[2pt]
K&=(-3;2)
\end{aligned}
\]
\end{zadaniReseniBox}
\begin{zadaniReseniBox}[title={Samostatná práce 2}]
$- x^{2} + 4x - 3\le0$
\tcblower
\[
\begin{aligned}
f(x)&=- x^{2} + 4x - 3\\[2pt]
f(x)=0&\ \Rightarrow\ K_0=\{1;3\}\\[2pt]
K&=(-\infty;1\rangle\cup\langle3;+\infty)
\end{aligned}
\]
\end{zadaniReseniBox}
\blok{Volitelná rezerva}
\poznamka{Studenti si vymyslí obdobný příklad se změnou jednoho koeficientu a vysvětlí, jak se změní počet kořenů či řešení.}
\blok{Domácí úkol}
\begin{zadaniReseniBox}[title={Domácí úkol}]
\begin{enumerate}[label=\alph*),itemsep=2pt]
\item $x^{2} - 7x + 12<0$
\item $- x^{2} + x + 12\ge0$
\item $2x^{2} - 8x + 6>0$
\item $x^{2} - 2x - 8\le0$
\item $x^{2} - 6x + 9\ge0$
\item $- 2x^{2} - 2x + 4<0$
\end{enumerate}
\tcblower
\begin{enumerate}[label=\alph*),itemsep=2pt]
\item
\[\begin{aligned}
f(x)&=x^{2} - 7x + 12\\[2pt]
f(x)=0&\ \Rightarrow\ K_0=\{3;4\}\\[2pt]
K&=(3;4)
\end{aligned}\]
\item
\[\begin{aligned}
f(x)&=- x^{2} + x + 12\\[2pt]
f(x)=0&\ \Rightarrow\ K_0=\{-3;4\}\\[2pt]
K&=\langle-3;4\rangle
\end{aligned}\]
\item
\[\begin{aligned}
f(x)&=2x^{2} - 8x + 6\\[2pt]
f(x)=0&\ \Rightarrow\ K_0=\{1;3\}\\[2pt]
K&=(-\infty;1)\cup(3;+\infty)
\end{aligned}\]
\item
\[\begin{aligned}
f(x)&=x^{2} - 2x - 8\\[2pt]
f(x)=0&\ \Rightarrow\ K_0=\{-2;4\}\\[2pt]
K&=\langle-2;4\rangle
\end{aligned}\]
\item
\[\begin{aligned}
f(x)&=x^{2} - 6x + 9\\[2pt]
f(x)=0&\ \Rightarrow\ K_0=\{3\}\\[2pt]
K&=\mathbb R
\end{aligned}\]
\item
\[\begin{aligned}
f(x)&=- 2x^{2} - 2x + 4\\[2pt]
f(x)=0&\ \Rightarrow\ K_0=\{-2;1\}\\[2pt]
K&=(-\infty;-2)\cup(1;+\infty)
\end{aligned}\]
\end{enumerate}
\end{zadaniReseniBox}
\poznamka{Orientační čas: úvod 5 min, společné příklady 18 min, samostatně 15 min, kontrola a shrnutí 7 min. Rezerva je dobrovolná. Číselné parametry modelových aplikací nejsou skutečně naměřená data.}
\end{document}
